\section{来源审计与阅读地图：这堂课到底在回答什么}

本讲录制于 2023 年 1 月 24 日，Jason Wei 用一条连续的教学链连接两个当时刚形成的研究主题：第一，为什么 language-model loss 可以随规模平滑下降，而某些任务能力却像跨过阈值一样突然出现；第二，为什么 chain-of-thought（CoT，思维链）并不是对所有模型都有效，而要等模型具有足够基础能力后才产生显著增益。重写时以官方 37 页 Google Slides 为视觉主轴，其中第 1--36 页含教学内容，第 37 页仅为致谢；现场 OpenAI Playground 的无 CoT / 有 CoT 对照则从官方录像精确抽帧补回。

\lecturefigure{slide-01.jpg}{课程主题：Scaling unlocks emergent abilities in language models}{官方课件第 1 页}

标题中的 ``unlock'' 很重要。它没有声称参数量本身创造了一个神秘模块，而是说：当模型、数据、训练计算和使用方法共同跨入某个区域后，原先在任务指标上看不见的能力开始可测。读者应把本讲当成一门“如何解释曲线”的课，而不是把 ``emergence'' 当作营销词。

\lecturefigure{slide-02.jpg}{课程路线：涌现能力、U-shaped scaling、CoT、BIG-Bench、跨语言推理与 self-consistency}{官方课件第 2 页}

整套课件先讨论预训练规模与任务能力，再转向 prompting-time 的推理干预。两部分并非并列综述：CoT 恰好构成一个关键案例，说明一种方法在小模型上可能无效甚至有害，在大模型上却打开新的任务区域；self-consistency 又进一步说明，推理时增加采样预算也存在规模门槛。

\begin{importantbox}{本讲的核心问题}
给定一条模型规模--性能曲线，我们怎样区分平滑改进、测量阈值、真正的机制变化和提示方法带来的增益？更进一步，当数据、后训练、metric 与 inference compute 都会移动曲线时，应该把哪一个量放在横轴上？
\end{importantbox}

\teachervoice{讲者开场强调，这项工作的重要性之一，是 Google、DeepMind 与 Stanford 的研究者形成了共同的观察框架。它首先是一套让不同模型、任务和团队可以对话的语言，而不是关于某个固定参数阈值的自然定律。}

\subsection{证据边界与重写原则}

这份讲义严格恢复课堂日期前的论证。官方课程页推荐 Emergent Abilities、Chain-of-Thought Prompting 与 Scaling Instruction-Finetuned Language Models 三篇论文；课件还直接使用 scaling laws、BIG-Bench、BIG-Bench Hard、multilingual CoT 与 self-consistency 等一手结果。2023 年稍后出现的“涌现是否是 metric artifact”争论具有研究价值，但它不属于 1 月 24 日课堂内容，因此不会被倒灌成讲者当时的结论。

\begin{warningbox}{历史材料不能被后见之明改写}
本讲中的 ``emergent'' 是操作性术语：小模型在所选 metric 上接近随机，大模型明显超过随机。它不自动等价于内部出现了一个离散新电路，也不证明阈值与模型参数之间存在唯一因果关系。讲义会保留课堂提出的证据，同时明确 metric、数据和模型家族带来的限制。
\end{warningbox}

\subsection{本章小结}

阅读后续图表时，需要始终携带三个问题：横轴究竟是什么；纵轴是否有随机或人类 baseline；曲线的拐点来自模型能力、数据质量、后训练、prompt，还是 metric 的共同作用。后文所有“涌现”判断都应回到这三问。

